\documentclass[10pt,twocolumn]{article}
\usepackage[letterpaper,margin=0.72in,columnsep=0.24in]{geometry}
\usepackage{times}
\usepackage{helvet}
\usepackage{courier}
\usepackage{amsmath,amssymb,bm}
\usepackage{algorithm}
\usepackage{algorithmic}

\setlength{\parskip}{2pt}
\setlength{\parindent}{1em}

\newcommand{\method}{IG-DS-SSM}
\newcommand{\pfim}{PFIM}
\newcommand{\dg}{DG-DS-SSM}

\begin{document}

\section{Method}

\subsection{Overview}

Given an input SAR image, the detector extracts multi-scale features using a lightweight Mamba-YOLO backbone. For a feature map $\mathbf{F}\in\mathbb{R}^{H\times W\times C}$, we flatten it into a token sequence
\begin{equation}
  \mathbf{x}=\mathrm{Flatten}(\mathbf{F})\in\mathbb{R}^{N\times C}, \quad N=HW.
\end{equation}
The proposed IG-DS-SSM block transforms $\mathbf{x}$ through de-speckled input modeling, importance estimation, and selective state-space scanning:
\begin{equation}
  \mathbf{y}=\mathrm{IG\mbox{-}DS\mbox{-}SSM}(\mathbf{x}, H, W).
\end{equation}
The output is reshaped back to two-dimensional features and passed to the detector neck and head.

\subsection{De-Speckled Input Modeling}

SAR speckle appears as local fluctuations that may be confused with weak ship structures. We estimate a local background representation $\bar{\mathbf{x}}$ using a depth-wise local operator:
\begin{equation}
  \bar{\mathbf{x}} = \mathrm{DWConv}(\mathbf{x}; H,W),
\end{equation}
where the operator is applied on the two-dimensional feature layout. The local residual is
\begin{equation}
  \mathbf{r} = \mathbf{x}-\bar{\mathbf{x}}.
\end{equation}
We then construct a de-speckled input
\begin{equation}
  \widetilde{\mathbf{x}} = \mathbf{x} - \lambda_s \mathbf{r},
  \label{eq:debiased-twocol}
\end{equation}
where $\lambda_s$ controls the strength of residual suppression. This operation does not attempt to remove all high-frequency information. Instead, it produces a state-space input whose local disturbance component is attenuated before recurrent propagation.

\subsection{Energy-Aware State Damping}

The local background energy provides a coarse reliability cue. For token $t$, we compute
\begin{equation}
  e_t = \frac{\|\bar{\mathbf{x}}_t\|_2}{\sqrt{C}}.
\end{equation}
The energy term is converted into a soft retention gate
\begin{equation}
  g^{E}_t = g_{\min} + (1-g_{\min})\sigma(-\alpha e_t),
\end{equation}
where $g_{\min}$ prevents complete state erasure and $\alpha$ controls damping strength. High-energy unstable background regions receive stronger damping, reducing the chance that clutter-dominated historical states persist across the scan.

\subsection{Perceptual Feature Importance Module}

Energy alone is insufficient because ship targets may also produce high responses. PFIM estimates a token-level importance score to distinguish target-likely regions from noise-likely regions. Given the de-speckled input, local background, and local statistics, PFIM outputs
\begin{equation}
  I_t = \mathrm{PFIM}(\widetilde{\mathbf{x}}_t,\bar{\mathbf{x}}_t,\mathrm{Std}_{local}(\mathbf{x})_t),
\end{equation}
where $I_t\in[0,1]$. In practice, PFIM is a lightweight branch using local two-dimensional statistics and point-wise transformations, followed by a sigmoid activation. The centered importance is denoted as
\begin{equation}
  \hat{I}_t = I_t - 0.5.
\end{equation}

\subsection{Importance-Guided State Update}

The selective state-space branch generates the input-dependent transition and projection terms from $\widetilde{\mathbf{x}}$. We use PFIM to modulate three update paths.

\paragraph{Memory retention.}
The importance-aware retention gate is
\begin{equation}
  g_t = g_{\min} + (1-g_{\min})\sigma(-\alpha e_t + \beta \hat{I}_t).
  \label{eq:gamma-twocol}
\end{equation}
This gate retains more state information around target-likely tokens and suppresses noise-likely propagation.

\paragraph{State writing.}
The write gate controls how strongly the current input enters the state:
\begin{equation}
  w_t = 1 + \rho \tanh(\hat{I}_t),
\end{equation}
\begin{equation}
  \mathbf{u}^{write}_t = w_t \odot \mathbf{u}_t.
\end{equation}

\paragraph{Dynamic transition step.}
PFIM also modulates the transition step:
\begin{equation}
  \Delta_t = \mathrm{softplus}(\Delta^{0}_t + \eta I_t),
\end{equation}
where $\Delta^{0}_t$ is the original transition logit. This allows the model to adjust state evolution speed according to target importance.

The resulting recurrent update can be written conceptually as
\begin{equation}
  \mathbf{h}_t =
  g_t \odot \bar{\mathbf{A}}_t\mathbf{h}_{t-1}
  + \bar{\mathbf{B}}_t\mathbf{u}^{write}_t,
\end{equation}
\begin{equation}
  \mathbf{y}_t = \mathbf{C}_t\mathbf{h}_t + \mathbf{D}\widetilde{\mathbf{x}}_t.
\end{equation}
The residual $\mathbf{D}\widetilde{\mathbf{x}}_t$ preserves current de-speckled evidence outside the recurrent state.

\begin{algorithm}[t]
\caption{IG-DS-SSM Forward Pass}
\begin{algorithmic}[1]
\REQUIRE Tokens $\mathbf{x}\in\mathbb{R}^{N\times C}$ and spatial size $(H,W)$
\STATE Estimate local background $\bar{\mathbf{x}}=\mathrm{DWConv}(\mathbf{x};H,W)$
\STATE Compute residual $\mathbf{r}=\mathbf{x}-\bar{\mathbf{x}}$
\STATE Build de-speckled input $\widetilde{\mathbf{x}}=\mathbf{x}-\lambda_s\mathbf{r}$
\STATE Estimate importance $\mathbf{I}=\mathrm{PFIM}(\widetilde{\mathbf{x}},\bar{\mathbf{x}},\mathrm{Std}_{local})$
\STATE Compute retention gate $g_t$, write gate $w_t$, and transition step $\Delta_t$
\STATE Run selective scan using gated state retention and gated input writing
\STATE Return output tokens $\mathbf{y}$
\end{algorithmic}
\end{algorithm}

\subsection{Decoupled-Gated Extension}

We further study DG-DS-SSM as an extension that explicitly separates erase, write, and transition controls. Its retain/erase branch is
\begin{equation}
  r_t = r_{\min} + (1-r_{\min})\sigma(-\alpha e_t+\beta_e\hat{I}_t),
\end{equation}
the write branch is
\begin{equation}
  q_t = 1+s_w\tanh(\rho_w\hat{I}_t),
\end{equation}
and the delta branch modulates $\Delta_t$ independently. DG-DS-SSM is useful for analysis because it clarifies whether old-state suppression, current-input writing, or transition adaptation contributes most to SAR detection. In our current results, it is competitive but does not surpass the full PFIM-enhanced DS-SSM line; therefore, we treat it as an interpretable extension rather than the main model.

\end{document}
